\documentclass[12pt]{article}
\usepackage[a4paper,margin=1in]{geometry}
\usepackage{amsmath,amssymb}
\usepackage{enumitem}
\usepackage{mathtools}
\usepackage{hyperref}
\usepackage{titlesec}
\usepackage{amsthm}

\title{Assignment 1, Matrix Theory}
\author{Simon Gustafsson}
\date{27 September 2026}

\newtheorem{theorem}{Theorem}
\newtheorem{lemma}{Lemma}
\newtheorem{corollary}{Corollary}
\newtheorem{definition}{Definition}
\newtheorem{remark}{Remark}


\begin{document}
\maketitle

\section*{Exercise 1}

\begin{proof}
Let $x, y \in \mathcal{C}$ and $z = \lambda x + (1-\lambda)y$ for some $\lambda \in (0,1)$.

\medskip
\noindent $(\Rightarrow)$ Assume $a_i^T z = b_i$. By linearity,
\[
a_i^T z = a_i^T(\lambda x + (1-\lambda)y) = \lambda\, a_i^T x + (1-\lambda)\, a_i^T y = b_i.
\]
Since $x, y \in \mathcal{C}$, $a_i^T x \leq b_i$ and $a_i^T y \leq b_i$. 
Suppose for contradiction that $a_i^T x < b_i$. Then, using $\lambda,(1 - \lambda) \in (0, 1)$,
\[
\lambda\, a_i^T x + (1-\lambda)\, a_i^T y < \lambda\, b_i + (1-\lambda)\, b_i = b_i,
\]
contradicting $a_i^T z = b_i$. Hence $a_i^T x = b_i$, with the same reasoning $a_i^T y = b_i$.

\medskip
\noindent $(\Leftarrow)$ Assume $a_i^T x = b_i$ and $a_i^T y = b_i$. Then $a_i^T z = \lambda\, a_i^T x + (1-\lambda)\, a_i^T y = \lambda\, b_i + (1-\lambda)\, b_i = b_i$.

\end{proof}

\end{document}
